\documentclass[letterstyle,12pt]{extarticle}
\usepackage[utf8]{inputenc}
\usepackage[nomarginpar, margin=0.5in]{geometry}
\usepackage{tabularx}
\usepackage[document]{ragged2e}
\usepackage{setspace}
\usepackage{amsmath}
\usepackage{amssymb}
\usepackage{xcolor}
\usepackage{mlmodern}
\usepackage{float}
\usepackage{supertabular}
% \usepackage[T1]{fontenc}
\usepackage{mathptmx}
\linespread{1.3} % Set line spacing 
\usepackage{pgfplots}
\usepackage{bigstrut}
\usepackage[ngerman]{babel}
\usepackage{mathtools}
\usepackage[T1]{fontenc}
\usepackage{enumitem}
\newcounter{qcounter}
\usepackage{multicol}
\pgfplotsset{compat = newest}
\usepackage{fancyhdr}
\everymath{\displaystyle}
\definecolor{agua}{RGB}{90,160,254}
\newcommand{\nspace}{\hspace{2.52mm}}

\newcommand\Tstrut{\smash[b]{\strut}}
\newcommand\Bstrut{\smash[t]{\strut}}
\newcommand{\dd}{\text{d}}
\newcommand{\di}{\, \text{d}}

\newcommand{\oldchoices}[1]{
\vspace{0.8em} 
\begin{enumerate}[label=(\Alph*)]
\setlength\itemsep{1em} #1 \interlinepenalty10000 
\end{enumerate}
}

\newcommand{\choices}[5]{
\vspace{0.8em} 
\begin{enumerate}[label=(\Alph*)]
\setlength\itemsep{1em} 
\item
#1 
\item 
#2
\item
#3
\item
#4
\item
#5
\end{enumerate}
}


\newcommand{\choicesline}[5]{    
\vspace{2em} \break 
\begin{tabularx}{0.95 \textwidth} { 
>{\arraybackslash}X 
>{\arraybackslash}X 
>{\arraybackslash}X 
>{\arraybackslash}X 
>{\arraybackslash}X }
(A) \; #1
& 
(B) \; #2
& 
(C) \; #3
& 
(D) \; #4
&
(E) \; #5 
\end{tabularx}
\vspace{2em} \break
}

\pgfplotsset{every axis/.append style={
    axis on top = true,
    width = \textwidth,  
    height =  \textwidth,
    clip=false, 
    axis line style = ultra thick, 
    axis lines = middle,
    label style={font=\Huge}, %x- and y-label sizes
    x label style={at={(axis description cs:1.04,0.03)},anchor=north},
    y label style={at={(axis description cs:-0.135,1.13)}}
    every tick label/.append style={font=\Huge},
    every major tick/.append style={very thick, major tick length=10pt},
    tick style = {xticklabel style={yshift=-0.2em, anchor=south}}
    }}
\pgfplotsset{tick style = {xticklabel style={yshift=-4ex, anchor=south}}}
\newcommand{\drawOrigin}{\draw (0, 0) node[left, below, xshift = -0.9em] {\Large $O$}} % draws the "O" marking the origin

\newcommand{\plotCircle}[4]{\draw #1 circle (#4 pt) [#2, fill = #3]}

\newcommand{\qspace}[1]{\begin{spacing}{1.8}#1\end{spacing}}
\newcommand{\ans}[1]{{\color{black} #1}}

\newenvironment{question}
    {\begin{minipage}{0.9 \textwidth}
        \item
    }
    { 
    \end{minipage} \vspace{4ex}
    }
    %

\newcommand{\bq}{\begin{question}}
\newcommand{\eq}{\end{question}}

\newcommand{\ques}[1]{\begin{minipage}{0.9 \textwidth} \item #1 \end{minipage} \vspace{4ex}}
\newcommand{\romanlist}[1]{\begin{enumerate}[label=\Roman*., leftmargin=15mm] #1 \end{enumerate}}

   
\pagestyle{fancy}
\lhead{\begin{figure}[H] \vspace{10ex} \includegraphics[scale=0.1]{logo-light.png} \vspace{-7.3ex} \end{figure}}
\rhead{\large \textsc{Review Exercises—Limits and Continuity}}
\fancyfoot{}
\lfoot{\copyright \textsc{Valculus} 2023}
\cfoot{\textbf{-\thepage-}}

\geometry{
    head = 1.1in,
    bottom = 1.3in,
    left = 0.5in,
    right = 0.5in
    }
    
\renewcommand{\headrulewidth}{0.3pt}
\renewcommand{\footrulewidth}{0.4pt}
\setlength{\headsep}{0.9in}

\begin{document}
\begin{center}
    {\bf
    \MakeUppercase{Limits and Continuity\\[1ex]
    NUMBER OF QUESTIONS—45 \\[1ex] 
    NO CALCULATOR \\[5ex]
    }}
\end{center}

\begin{list}{\textbf{\arabic{qcounter}.}~}{\usecounter{qcounter}}
\setlength\itemsep{3em}

\begin{question}
If \(\lim_{x \to 3} f(x) = 8,\) then which option is true?
\choices
{\(f(3) = 8.\)}
{\ans{As \(x\) approaches \(3,\) \(f(x)\) approaches \(8.\)}}
{\(f(8) = 3.\)} 
{As \(x\) approaches \(8,\) \(f(x)\) approaches \(3.\)} 
{\(f\) is continuous at \(x = 3.\)}
\end{question}

\begin{question}
\(\lim_{x \to 0} \frac{\sin 3x}{x}\) is 
\choicesline
{\(0\)}
{\(\frac{1}{3}\)}
{1}
{\ans 3}
{nonexistent}
\end{question}

\begin{question}
\(\lim_{x \to -7} \frac{x + 7}{49 - x^2}\) is 
\choicesline
{\(-7\)}
{\ans{\(-\frac{1}{14}\)}}
{\(-\frac{1}{49}\)}
{\(\frac{1}{14}\)}
{\(49\)}
\end{question}

\begin{question}
\(\lim_{x \to 0} \frac{x}{\tan x}\) is 
\choicesline
{\(-1\)}
{0}
{\ans 1}
{\(\pi\)}
{nonexistent}
\end{question}

\begin{question}
\(\lim_{x \to \infty} \frac{3x^4 - 8x^3 + x^2 - 10}{2x^2 - 5x^3 - 10x^4}\) is 
\choicesline
{\ans{\(-\frac{3}{10}\)}}
{\(0\)}
{\(\frac{3}{10}\)}
{\(\frac{3}{2}\)}
{nonexistent}
\end{question}

\begin{question}
The horizontal asymptote of \(f(x) = \frac{3x^3 + x^2 - 4}{8 - 5x^3}\) \, is \\
\choicesline
{\(y = -\frac{5}{3}\)}
{\ans{\(y = -\frac{3}{5}\)}}
{\(y = \frac{3}{8}\)}
{\(y = \frac{3}{5}\)}
{\(y = \frac{5}{3}\)}
\end{question}

\begin{question}
\(\lim_{x \to \pi/2} \tan 2x\) is
\choicesline
{\(-1\)}
{\(-\frac{1}{2}\)}
{\ans 0}
{1}
{\(\pi\)}
\end{question}

\begin{question}
Function \(g\) is discontinuous at \(x = 5.\)
Selected values of $g$ are shown in the table below. \\[1ex]
\renewcommand{\arraystretch}{1.2}
\begin{center}
\begin{tabularx}{0.7\textwidth} { 
  | >{\arraybackslash}X 
  || >{\centering\arraybackslash}X 
  | >{\centering\arraybackslash}X 
  | >{\centering\arraybackslash}X 
  | >{\centering\arraybackslash}X 
  | >{\centering\arraybackslash}X | }
 \hline
 $\hspace{2.2ex} x$ & 4.99 & 4.999 & 5 & 5.0001 &5.001  \\
 \hline
 $f(x)$ & 2.99 & 2.999 & $-4$ & 3.001  &3.01 \\
 \hline
\end{tabularx} \\[1.5ex]
\end{center}
A reasonable estimate for \(\lim_{x \to 5} g(x)\) is \\
\choicesline
{\(-4\)}
{\(-3\)}
{2}
{\ans 3}
{5}
\end{question}

\begin{question}
If
\(f(x)= \begin{cases} 
3x - 2 \cos x& x < \pi \\
x^2 & x \geqslant \pi, \\ 
\end{cases} 
\) \, 
then \(\lim_{x \to \pi^-} f(x)\) is \\
\choicesline
{\(-\pi^2\)}
{\(3 \pi\)}
{\(3 \pi - 2\)}
{\ans{\(3 \pi + 2\)}}
{\(\pi^2\)}
\end{question}

\begin{question}
\(\lim_{x \to 3} \frac{\sqrt{x - 2} - 1}{9 - 3x}\) is 
\choicesline
{\(-\frac{1}{2}\)}
{\ans{\(-\frac{1}{6}\)}}
{\(\frac{1}{6}\)}
{\(\frac{1}{2}\)}
{nonexistent}
\end{question}

\begin{question}
Given that \(\lim_{x \to a} f(x)\) exists, which statements must be true?
\romanlist{
\item
\(f(x)\) is continuous at \(x = a.\)
\item
\(\lim_{x \to a^-} f(x) = \lim_{x \to a^+} f(x).\)
\item
\(f(a)\) is defined.
}
\choices
{I only}
{\ans{II only}}
{I and II only}
{II and III only}
{I, II, and III}
\end{question}

\begin{question}
If \(\lim_{k \to 0} \frac{e^k - 1}{k} = 1,\) then \(\lim_{k \to 0} \frac{e^2 - e^{2 + k}}{k}\) is \\
\choicesline
{\ans{\(-e^2\)}}
{\(-e^{-2}\)}
{\(e^{-2}\)}
{1}
{\(e^2\)}
\end{question}

\begin{question}
\(g(x) = \frac{x^2 - 5x + 6}{x - 3}\) \; has a removable discontinuity at \\ 
\choicesline
{\(x = -3\)}
{\(x = -2\)}
{\(x = 0\)}
{\(x = 2\)}
{\ans{\(x = 3\)}}
\end{question}

\begin{question}
\(\lim_{x \to \infty} \frac{\sin(2x)}{x + 2}\) is 
\choicesline
{\(-1\)}
{\ans 0}
{1}
{2}
{\(\infty\)}
\end{question}

\begin{question}
Let
\( f(x) = \begin{cases} 
2 - kx & x\leqslant 3 \\
kx^2 -22 & x > 3. \\ 
\end{cases}
\)
For what value of \(k\) is \(f\) continuous at \(x = 3 \; ?\) \\
\choicesline
{\(-4\)}
{0}
{1}
{\ans 2}
{3}
\end{question}

\begin{question}
The oblique asymptote of \(f(x) = \frac{x^2 + 7x + 1}{x - 2}\) is \\ 
\choicesline
{\(y = 9\)}
{\(y = x - 9\)}
{\(y = x\)}
{\(y = x + 5\)}
{\ans{\(y = x + 9\)}}
\end{question}

\begin{question}
\(\lim_{x \to \infty} \frac{\sqrt{4x^2 - 1}}{x + 3} \, \) is 
\choicesline
{\(-4\)}
{\(-2\)}
{\ans 2}
{4}
{nonexistent}
\end{question}

\begin{minipage}{\textwidth}
Questions \ref{q:graph-start}--\ref{q:graph-end} refer to the following graph.
%
\begin{center}
\begin{tikzpicture}
\begin{axis}[
    axis lines = middle,
    xlabel=\large $x$,
    ylabel=\large $y$,
    every axis x label/.style={at={(current axis.right of origin)},anchor=west},
    every axis y label/.style={at=(current axis.above origin),anchor=south},
    xmin = -5.5, xmax = 5.5,
    ymin = -5.5, ymax = 5.5,
    xticklabel={
        % test if the x value is below zero
        \ifdim \tick pt < 0pt
          % if yes, calculate the absolute value
          \pgfmathparse{abs(\tick)}
          % and print first a minus sign in a zero-width box, followed by the absolute value
          \llap{$-{}$}\pgfmathprintnumber{\pgfmathresult}
      \else
         % if no, print as usual
          \pgfmathprintnumber{\tick}
      \fi
        },
    xtick distance = 1,
    ytick distance = 1,
    grid = both,
    major grid style = {lightgray},
    minor grid style = {lightgray!25},
    width = 0.7 \textwidth,
    height = 0.7 \textwidth]
    \addplot[red, ultra thick, smooth, domain=-5:0]{x + 4};
    \addplot[red, ultra thick, smooth, domain=0:4.8]{2 - x}node[pos=1.07]{\Large $f$};
    \addplot[agua!70, ultra thick, smooth, domain=-4.8:0]{(x + 2)^2 - 3};
    \addplot[agua!70, ultra thick, smooth, domain=0:4.8]{2 + (1/2)* sin(deg(x))}node[pos=1.07]{\Large $g$};
    \plotCircle{(0, 2)}{red}{red}{3};
    \plotCircle{(0, 1)}{agua!70}{agua!70}{3};
    \plotCircle{(0, 4)}{red}{white}{3};
\end{axis}
\end{tikzpicture}
\end{center}
\end{minipage}

\begin{question}
\label{q:graph-start}
\(\lim_{x \to 0} f(x)\) is 
\choicesline
{0}
{1}
{2}
{4}
{\ans{nonexistent}}
\end{question}

\begin{question}
\(\lim_{x \to 0^+} g(x)\) is
\choicesline
{0}
{1}
{\ans 2}
{4}
{nonexistent}
\end{question}

\begin{question}
\(\lim_{x \to -2} [f(x) - g(x)]\) is \\
\choicesline
{\(-5\)}
{\(-1\)}
{1}
{\ans 5}
{nonexistent}
\end{question}

\begin{question}
\(\lim_{x \to -2} [g(x)]^2\) is 
\choicesline
{\(-4\)}
{\(-2\)}
{2}
{4}
{\ans 9}
\end{question}

\begin{question}
\(\lim_{x \to 0} [f(x) g(x)]\) is 
\choicesline
{1}
{2}
{\ans 4}
{8}
{nonexistent}
\end{question}

\begin{question}
\label{q:graph-end}
\(\lim_{x \to 0} f(-x^2)\) is 
\choicesline
{0}
{1}
{2}
{\ans 4}
{nonexistent}
\end{question}

\begin{question}
On what interval of \(x\) is \(f(x) = \frac{\ln(x - 2)}{x^2 - 9}\) continuous?
\choices
{\((-\infty, -3)\cup(3, \infty)\)}
{\((-\infty, -2)\)}
{\((2, \infty)\)}
{\((3, \infty)\)}
{\ans{\((2, 3)\cup(3, \infty)\)}}
\end{question}

\begin{question}
\(\lim_{x \to -3} \frac{5\sin(x + 3)}{6 + 2x}\) is 
\choicesline
{0}
{\(\frac{5}{6}\)}
{1}
{\ans{\(\frac{5}{2}\)}}
{nonexistent}
\end{question}

\begin{question}
Function \(f\) is continuous and satisfies \(f(4) = 8.\) If \(\lim_{x \to 2} g(x) = 4,\) then \(\lim_{x \to 2} f(g(x))\) is \\
\choicesline
{\(-8\)}
{\(-4\)}
{2}
{4}
{\ans 8}
\end{question}

\begin{question}
\(\lim_{x \to 2} \cfrac{\cfrac{1}{2 \Bstrut} - \cfrac{1}{x \Bstrut} }{\Tstrut 2 - x}\) is
\choicesline
{\(-\frac{1}{2}\)}
{\ans{\(-\frac{1}{4}\)}}
{\(\frac{1}{4}\)}
{\(\frac{1}{2}\)}
{nonexistent}
\end{question}

\begin{question}
\(\lim_{x \to 0} x \sin\left( \frac{1}{x^2}\right)\) is 
\choicesline
{\ans 0}
{\(\frac{1}{4}\)}
{\(\frac{1}{2}\)}
{1}
{nonexistent}
\end{question}

\begin{question}
Let
\(f(x)= \begin{cases} 
x^2 + 1& x < p \\
2x & x \geqslant p. \\ 
\end{cases} 
\)
If \(f(x)\) is continuous at \(x = p,\) then \(p\) is \\
\choicesline
{\ans{\(-1\)}}
{0}
{1}
{2}
{4}
\end{question}

\begin{question}
\(\lim_{t \to \infty} \sin t\) is 
\choicesline
{\(-1\)}
{0}
{1}
{\(\pi\)}
{\ans{nonexistent}}
\end{question}

\begin{question}If \(\lim_{x \to 2} f(x) = -3,\)
then \(\lim_{x \to 2} \left([f(x)]^2 - 2x \right)\) is \\
\choicesline
{\(-7\)}
{\(-3\)}
{4}
{\ans 5}
{9}
\end{question}

\begin{question}
Function \(f\) is continuous. Selected values of \(f(x)\) are shown in the table below. \\
\begin{center}
\begin{tabularx}{0.7\textwidth} { 
  | >{\arraybackslash}X 
  || >{\centering\arraybackslash}X 
  | >{\centering\arraybackslash}X 
  | >{\centering\arraybackslash}X 
  | >{\centering\arraybackslash}X 
  | >{\centering\arraybackslash}X | }
 \hline
 \(\hspace{2.2ex} x\) & \(-1\) & 2 & 3 & 6 & 11  \\
 \hline
 \(f(x)\) & 2 & 1 & \(1\) & 1  & 2 \\
 \hline
\end{tabularx} \\
\end{center}
Following the Intermediate Value Theorem, which value of \(f(x)\) is guaranteed to exist for \(-1 \leqslant x \leqslant 11 \; ?\) \\
\choicesline
{\ans 0}
{3}
{5}
{6}
{11}
\end{question}

\begin{question}
\(\lim_{x \to -\infty} \frac{\sqrt{4x^6 - 4x^2 + 1}}{3x^3 + 2} \, \) is 
\choicesline
{\ans{\(-\infty\)}}
{\(-\frac{2}{3}\)}
{\(\frac{2}{3}\)}
{\(\infty\)}
{nonexistent}
\end{question}

\begin{question}
\(\lim_{x \to \infty} x \sin \left(\frac{1}{x}\right)\) is 
\choicesline
{\(-1\)}
{0}
{\ans 1}
{\(\pi\)}
{nonexistent}
\end{question}

\begin{question}
At \(x = 4,\) which choice about \(g(x) = \frac{12 + x - x^2}{x - 4}\) is true?
\choices
{\(g(x)\) has a vertical asymptote at \(x = 4.\)}
{\(g(x)\) has a jump discontinuity at \(x = 4.\)}
{\ans{\(g(x)\) has a removeable discontinuity at \(x = 4.\)}}
{\(\lim_{x \to 4} g(x)\) does not exist.}
{\(g(x)\) is continuous at \(x = 4.\)}
\end{question}

\begin{question}
\(\lim_{x \to 0^+} \ln(\sin x)\) is 
\choicesline
{\ans{\(-\infty\)}}
{0}
{1}
{\(e\)}
{\(\infty\)}
\end{question}

\begin{question}
\(\lim_{x \to 0} \frac{x^2}{\sin^2 2x}\) is 
\choicesline
{0}
{\ans{\(\frac{1}{4}\)}}
{1}
{4}
{nonexistent}
\end{question}

\begin{question}
\(\lim_{x \to 0} \frac{x - x\cos(x)}{x^2}\) is 
\choicesline
{\(-\pi\)}
{\ans 0}
{1}
{\(\pi\)}
{nonexistent}
\end{question}

\begin{question}
\(\lim_{x \to 0^+} \arctan \left(\frac{1}{x}\right)\) is 
\choicesline
{\(-\infty\)}
{\(-\frac{\pi}{2}\)}
{0}
{\ans{\(\frac{\pi}{2}\)}}
{\(\infty\)}
\end{question}

\begin{question}
\(\lim_{x \to \infty} \frac{5e^{x} - x}{8e^{x} + 9}\) is 
\choicesline
{\(-\frac{5}{8}\)}
{0}
{\ans{\(\frac{5}{8}\)}}
{1}
{\(\infty\)}
\end{question}

\begin{question}
\(\lim_{x \to \pi/4} \, \frac{\cos 2x}{\cos x - \sin x}\) is 
\choicesline
{0}
{\(\frac{\sqrt 2}{2}\)}
{\ans{\(\sqrt 2\)}}
{\(\pi\)}
{nonexistent}
\end{question}

\begin{question}
Functions \(g\) and \(h\) are continuous and satisfy \(g(1) = h(1) = 3.\) Function \(f\) satisfies \(g(x) \leqslant f(x) \leqslant h(x)\) for \(0 \leqslant x \leqslant 2.\) 
Which statements must be true?
\romanlist{
\item
\(\lim_{x \to 1} g(x) = \lim_{x \to 1} h(x) = 3.\) 
\item
\(\lim_{x \to 1} f(x) = 3.\)
\item
\(f(x)\) is continuous at \(x = 1.\)
}
\choices
{I only}
{II only}
{I and II only}
{II and III only}
{\ans{I, II, and III}}
\end{question}

\begin{question}
If \(f(1) = 2,\) \(\lim_{x \to 1^-} f(x) = 4,\) and \(\lim_{x \to 1^+} f(x) = -1,\) then \(\lim_{x \to 1} f(\cos(x - 1))\) is 
\choicesline
{\(-1\)}
{0}
{2}
{\ans 4}
{nonexistent}
\end{question}

\begin{question}
If \(\lim_{x \to 3} (2x + 4) = 10,\) then \(|(2x + 4) - 10| < \epsilon\) and \(|x - 3| < \delta,\) where \(\delta =\) \\
\choicesline
{\(\frac{\epsilon}{4}\)}
{\ans{\(\frac{\epsilon}{2}\)}}
{\(\epsilon\)}
{\(2 \epsilon\)}
{\(4 \epsilon\)}
\end{question}

\begin{question}
If \(\lim_{x \to a} f(x) = \infty,\) then which option is true?
\choices
{For positive \(M,\) there exists a positive \(\delta\) such that \(f(x) > M\) for \(|x - a| > 0.\)}
{For positive \(M,\) there exists a positive \(\delta\) such that \(f(x) < M\) for \(|x - a| > \delta.\)}
{For positive \(M,\) there exists a positive \(\delta\) such that \(f(x) > M\) for \(|x - a| > \delta.\)}
{For positive \(M,\) there exists a positive \(\delta\) such that \(f(x) < M\) for \(0 < |x - a| < \delta.\)}
{\ans{For positive \(M,\) there exists a positive \(\delta\) such that \(f(x) > M\) for \(0 < |x - a| < \delta.\)}}
\end{question}


\end{list}

\clearpage 

\twocolumn

\begin{supertabular}{ll}
% e.g., 1 & B \\
{\bf1}. & B \\ 
{\bf2}. & D \\ 
{\bf3}. & B \\ 
{\bf4}. & C \\ 
{\bf5}. & A \\ 
{\bf6}. & B \\ 
{\bf7}. & C \\ 
{\bf8}. & D \\ 
{\bf9}. & D \\ 
{\bf10}. & B \\ 
{\bf11}. & B \\ 
{\bf12}. & A \\ 
{\bf13}. & E \\ 
{\bf14}. & B \\ 
{\bf15}. & D \\ 
{\bf16}. & E \\ 
{\bf17}. & C \\ 
{\bf18}. & E \\ 
{\bf19}. & C \\ 
{\bf20}. & D \\ 
{\bf21}. & E \\ 
{\bf22}. & C \\ 
{\bf23}. & D \\ 
{\bf24}. & E \\ 
{\bf25}. & D \\ 
{\bf26}. & E \\ 
{\bf27}. & B \\ 
{\bf28}. & A \\ 
{\bf29}. & A \\ 
{\bf30}. & E \\ 
{\bf31}. & D \\ 
{\bf32}. & A \\ 
{\bf33}. & A \\ 
{\bf34}. & C \\ 
{\bf35}. & C \\ 
{\bf36}. & A \\ 
{\bf37}. & B \\ 
{\bf38}. & B \\ 
{\bf39}. & D \\ 
{\bf40}. & C \\ 
{\bf41}. & C \\ 
{\bf42}. & E \\ 
{\bf43}. & D \\ 
{\bf44}. & B \\ 
{\bf45}. & E \\ 

\end{supertabular}
\end{document}